\section{Whitebox PIT for \texorpdfstring{$\SPSU{[k]}{}{}$}{SkPSU}} \label{sec:pf:thm1}

We consider a bloated model of computation which naturally generalizes $\SPSU{}{}{}$ circuits and works ideally under the $\DiDI$-techniques.

\begin{definition}
We call a circuit $\calC \in \Gen{k}{s}$, over $\ringR(\x)$, for any ring $\ringR$, with parameter $k$ and size-$s$, if $\calC \in \Sigma^{[k]} (\PSU{}{}/\PSU{}{}) \cdot (\SWSU{}{}/\SWSU{}{})$. It computes $f \in \ringR(\x)$, if $f\;=\; \sum_{i=1}^k\, T_{i}$, where 
    \begin{itemize}
        \item $T_i\,=:\,(U_i/V_i)\cdot (P_i/Q_i)$, for $U_i, V_i \in \PSU{}{}$, and $P_i,Q_i \in \SWSU{}{}$,  
%
        \item $\size(T_i) = \size(U_i) + \size(V_i) + \size(P_i) + \size(Q_i)$, and $\size(f)=\sum_{i \in [k]} \size(T_i)$.
    \end{itemize}
\end{definition}

\noindent It is easy to see that all size-$s$ $\SPSU{[k]}{}{}$ circuit are in $\Gen{k}{s}$. We will design the \emph{recursive} algorithm on $\Gen{k}{s}$. 

\begin{proof}[Proof of \autoref{thm:thm1}] 
Begin with defining $T_{i,0}:=T_i$ and $f_{0}:=f$ where $T_{i,0} \in \PSU{}{}$;  $\sum_{i} T_{i,0}=f_0$, and $f_0$ has size $\le s$. Assume $\pdeg(f)<d\le s$; we keep the parameter $d$ separately, to help optimize the complexity later. In every recursive call we work with $\Gen{\cdot}{\cdot}$ circuits. 

As the input case, define $U_{i,0}:=T_{i,0}$ and $V_{i,0}:=P_{i,0}:= Q_{i,0}:=1$. We will use the hitting set of product of sparse polynomials (refer section \ref{sec:basic-act}) to obtain a point $\alphavec = (a_1, \dots, a_n) \in \F^n$ such that $U_{i,0} \rvert_{\x = \alphavec} \neq 0$, for all $i \in [k]$. Eventually this evaluation point will help in maintaining the invertibility of $\PSU{}{}$. Consider
\begin{align*}
    g \;:=\; \prod_{i \in [k]} T_{i,0} &\;=\; \prod_{i \in [k]} U_{i,0} \;=\; \prod_{i \in [\ell]} \sum_{j \in [n]} f_{ij}(x_j)\,,
\end{align*}
where $f_{ij}(x_j)$ are univariate polynomials of degree at most $d$ and $\ell \leq k \cdot s$. Note that $\deg g \leq d \cdot k \cdot s$ and $g$ is computable by a $\PSU{}{}$ circuit of size $O(s)$. Invoke \autoref{lem:prod-sparse-hs} to obtain a hitting set $\calH$, then evaluate $g$ on every point of $\calH$ to find an element $\alphavec \in \calH$ such that $g(\alphavec) \neq 0$. We emphasise that in whitebox setting all $U_{i,0}$, are readily available for evaluation. Since, the size of the set is $\poly(s)$ and each evaluation takes $\poly(s)$ time, this preliminary step will add $\poly(s)$ time to the overall time complexity. Moreover, we obtain the $\alphavec \in \F^n$ which possess the required property. 

To capture the non-zeroness, consider a 1-1 homomorphism $\Phi : \F[\x] \longrightarrow \F[\x,z]$ such that $x_i \mapsto z\cdot x_i + a_i$ where $a_i$ is the $i$-th coordinate of $\alphavec$, obtained earlier. Invertibility implies that $f_{0} = 0 \iff \Phi(f_{0}) = 0$. 
%
%
Now we proceed with the recursive algorithm which first reduces the identity testing from top-fanin $k$ to $k-1$. Note: $k=1$ is trivial.

\subsubsection*{First Step: Efficient reduction from $k$ to $k-1$}
By assumption, $\sum_{i=1}^k\,T_{i,0}\,=\,f_{0}$ and $T_{k,0}\ne0$. Apply $\Phi$ both sides, then divide and derive:
\begin{align} \label{Eq:thm1-k-k-1}
    \sum_{i \in [k]}\,T_{i,0} \;=\; f_0 \; &\iff\;  \sum_{i \in [k]}\,\Phi(T_{i,0}) \;=\; \Phi(f_0) \nonumber\\ 
    &\iff  \sum_{i \in [k-1]}\, \frac{\Phi(T_{i,0})}{\Phi(T_{k,0})}\,+\, 1 \;=\; \frac{\Phi(f_0)}{\Phi(T_{k,0})} \nonumber\\ 
    &\implies \sum_{i \in [k-1]}\,\partial_{z} \left( \frac{\Phi(T_{i,0})}{\Phi(T_{k,0})}\right) \;=\; \partial_{z}\left(\frac{\Phi(f_0)}{\Phi(T_{k,0})}\right) \nonumber\\ 
&\iff \sum_{i=1}^{k-1}\,\frac{\Phi(T_{i,0})}{\Phi(T_{k,0})} \cdot \dlg_{z} \left( \frac{\Phi(T_{i,0})}{\Phi(T_{k,0})} \right) \;=\; \partial_{z} \left( \frac{\Phi(f_0)}{\Phi(T_{k,0})}\right)\;.
\end{align}
Here onwards we say \(\dlg\) to mean \(\dlg_z\), unless stated otherwise. Define the following:
\begin{itemize}
    \setlength\itemsep{1em}
    \item $\ringR_1\,:=\, \F[z]/\langle z^{d} \rangle$. Note that, \autoref{Eq:thm1-k-k-1} holds over $\ringR_1(\x)$.
    \item $\widetilde{T}_{i,1}:=\Phi(T_{i,0})/\Phi(T_{k,0}) \cdot \dlg (\Phi(T_{i,0})/\Phi(T_{k,0}))$, $\forall\;i \in [k-1]$.
    \item $ f_{1}:=\partial_{z}(\Phi(f_0)/\Phi(T_{k,0}))$, over $\ringR_1(\x)$.
\end{itemize}
\vspace*{1em}

\textbf{Definability of $T_{i,1}$ and $f_1$.} It is easy to see that these are well-defined terms. Here, we emphasize that we do not exactly compute/store $\widetilde{T}_{i,1}$ as a fraction where the degree in $z$ is $<d$; instead it is computed as an element in $\F(z,\x)$, where $z$ is a formal variable. Formally, we compute $T_{i,1} \in \F(z,\x)$, such that $\widetilde{T}_{i,1} = T_{i,1}$, over $\ringR_1(\x)$. We keep track of the degree of $z$ in $T_{i,1}$. 
Thus, $\sum_{i \in [k-1]}\,T_{i,1}=f_1$, over $\ringR_1(\x)$.

\vspace*{1em}
\textbf{The `iff' condition.} To show that our one step of $\DiDI$ has reduced to the identity testing of $\Gen{k-1}{\cdot}$, we need an $\iff$ condition. So far equality in \autoref{Eq:thm1-k-k-1} over $\ringR_1(\x)$ is \emph{one-sided}. Note that $f_{1} \ne 0$ implies $\val_{z}(f_{1}) < d=:d_1$. By assumption, $\Phi(T_{k,0})$ is invertible over $\ringR_1(\x)$. Further, $f_{1} = 0$, over $\ringR_1(\x)$, which implies --

\begin{enumerate}
    \setlength\itemsep{1em}
    \item Either, $\Phi(f_0)/\Phi(T_{k,0})$ is $z$-free. Then $\Phi(f_0)/\Phi(T_{k,0}) \,\in \F(\x)$, which further implies it is in $\F$, because of the map $\Phi$ ($z$-free implies $\x$-free, by substituting $z=0$). Also, note that $f_0, T_{k,0} \ne 0$ implies $\Phi(f_0)/\Phi(T_{k,0})$ is a \emph{nonzero} element in $\F$. Thus, it suffices to check whether  $\Phi(f_0)\rvert_{z=0}$ is non-zero or not. 
%
    \item Or, $\partial_{z}(\Phi(f_0)/\Phi(T_{k,0}))= z^{d_1} \cdot p$ where $p \in \F(z,\x)$ s.t.~$\val_{z}(p) \ge 0$. By simple power series expansion, one can show that $p \in \F(x)[[z]]$. 
    
    \begin{lemma}[Valuation]\label{lem:val-powerseries}
        Consider $f \in \F(\x,y)$ such that $\val_{y}(f) \ge 0$. Then, $f \in \F(\x)[[y]]\,\bigcap\,\F(\x,y)$.
    \end{lemma}
    \begin{proof-sketch}
        Let $f=g/h$, where $g,h \in \F[\x,y]$. Now, $\val_{y}(f) \ge 0$, implies $\val_y(g) \ge \val_y(h)$. Let $\val_{y}(g)=d_1$ and $\val_{y}(h)=d_2$, where $d_1 \ge d_2 \ge 0$. Write $g=y^{d_1}\cdot \tilde{g}$ and $h=y^{d_2}\cdot \tilde{h}$. Write, $\tilde{h}= h_0+h_1\,y+h_2\,y^2+\hdots + h_d\,y^d$, for some $d$. Note that $h_0 \ne 0$. Thus, 
        \begin{align*}
            f&\;=\; y^{d_1-d_2} \cdot \tilde{g}/(h_0+h_1\,y+\hdots+h_d\,y^d) \\ &\;=\;y^{d_1-d_2} \cdot (\tilde{g}/h_0) \cdot (1+(h_1/h_0)\,y+\hdots + (h_d/h_0)\,y^d)^{-1}\; \in \F(\x)[[y]]\;.
        \end{align*}
        The last conclusion follows by the inverse identity in the power-series ring.
    \end{proof-sketch}

    Hence, $\Phi(f_0)/\Phi(T_{k,0}) = z^{d_1+1} \cdot q$ where $q \in F(\x)[[z]]$, i.e. \[\Phi(f_0)/\Phi(T_{k,0}) \in \langle z^{d_1+1}\rangle_{\F(\x)[[z]]} \implies \val_{z}(\Phi(f_0)) \ge d+1,\] a contradiction.
\end{enumerate}

\noindent Conversely, it is obvious that $f_0=0$ implies $f_1=0$. Thus, we have proved the following 
\[ \sum_{i \in [k]}\,T_{i,0}\,\ne 0 \;\text{ over}\;\F[\x] \iff \sum_{i \in [k-1]}\, T_{i,1} \ne 0 \;\text{ over} \;\ringR_1(\x),\;\text{ or}\,,\;\;0 \ne \Phi(f_0) \rvert_{z=0} \in \F\;.\]
 Eventually, we show that $T_{i,1} \in (\PSU{}{}/\PSU{}{}) \cdot (\SWSU{}{}/\SWSU{}{})$, over $\ringR_1(\x)$, with polynomial blowup in size (\autoref{cl:size-bound-on-Ti}). So, the above circuit is in $\Gen{k-1}{\cdot}$, over $\ringR_1(\x)$, which we recurse on to finally give the identity testing. The subsequent steps will be a bit more tricky:

\vspace*{1em}
\subsubsection*{Induction step} 

Assume that we are in the $j$-th step ($j\ge1$). Our induction hypothesis assumes --

\begin{enumerate}
    \item $\sum_{i \in [k-j]}\, T_{i,j} = f_j$, over $\ringR_j(\x)$, where $\ringR_j:= \F[z]/\langle z^{d_j}\rangle$ for \(d_j < d\), and $T_{i,j} \ne 0$. \label{hypo1}
    \item $\val_{z}(T_{i,j}) \ge 0, \forall i \in [k-j]$. \label{hypo2}
    \item Non-zero preserving iff condition
    \begin{align*} 
        f \neq 0 \text{, over } \F[\x] \iff & f_j \ne 0, \text{ over } \ringR_j(\x), \\
        & \text{or }\bigvee_{i=0}^{j-1} \left((f_i/T_{k-i,i}) \rvert_{z=0}\ne 0,\, \text{over}\,\F(\x)\right)
    \end{align*} \label{hypo3}
    \item Here, $T_{i,j}=: \left( U_{i,j}/V_{i,j} \right) \cdot \left( P_{i,j}/Q_{i,j} \right)$, where $U_{i,j}, V_{i,j} \in \PSU{}{}$, and $P_{i,j}, Q_{i,j} \in \SWSU{}{}$, each in $\ringR_j[\x]$. Think of them being computed as $\F(z,\x)$, with the degrees being tracked. Wlog, assume that $\val_{z}(T_{k-j,j})$ is the minimal among all $T_{i,j}$'s. \label{hypo4}
    \item $U_{i,j}\rvert_{z=0}, V_{i,j}\rvert_{z=0}  \in \F\backslash \{0\}$. \label{hypo5}
\end{enumerate} 

We follow as before without applying homomorphism any further. Note that the `or condition' in the hypothesis \ref{hypo3} is similar to the $j=0$ case except that there is no $\Phi$: this is because $\Phi(f_0)\rvert_{z=0} \ne 0 \iff \Phi(f_0/T_{k,0})\rvert_{z=0} \ne 0$. This condition just separates the derivative from the constant-term.

\textbf{Efficient reduction from $k-j$ to $k-j-1$.} Let $\val_{z}(T_{i,j})=:v_{i,j}$, for all $i \in [k-j]$. Note that \[\min_{i} \val_{z}(T_{i,j})=\min_{i}\val_{z}(P_{i,j}/Q_{i,j})=v_{k-j,j}\] since $\val_{z}(U_{i,j}) =\val_{z}(V_{i,j}) =0$ (else we reorder). We remark that  $ 0 \le v_{i,j} < d_j$ for all $i$'s in $j$-th step; upper-bound is strict, since otherwise $T_{i,j}=0$ over $\ringR_j(x)$.

Similar to the first step, we divide with $T_{k-j,j}$ which has $\min \val$ and then derive:
\begin{align}
    \sum_{i \in [k-j]}\,T_{i,j} \;=\; f_j  &\iff  \sum_{i \in [k-j-1]}\,T_{i,j}/T_{k-j,j} \,+\, 1 \;=\; f_j/T_{k-j,j} \nonumber\\  &\implies \sum_{i \in [k-j-1]}\,\partial_{z} (T_{i,j}/T_{k-j,j}) \;=\; \partial_{z}(f_j/T_{k-j,j})\nonumber \\ &\iff \sum_{i=1}^{k-j-1}\,T_{i,j}/T_{k-j,j} \cdot \dlg (T_{i,j}/T_{k-j,j}) \;=\; \partial_{z} (f_j/T_{k-j,j}) \label{eq:induction-thm1-main-eq}
\end{align}
Define the following:
\begin{itemize}
    \setlength\itemsep{1em}
    \item $\ringR_{j+1}:= \F[z]/ \langle z^{d_{j+1}} \rangle$, where $d_{j+1}:=d_j-v_{k-j,j}-1$. 
    \item $\widetilde{T}_{i,j+1} := T_{i,j}/T_{k-j,j} \cdot \dlg(T_{i,j}/T_{k-j,j})$, $\forall\;i \in [k-j-1]$.
    \item $f_{j+1}:=\partial_{z}(f_j/T_{k-j,j})$, over $\ringR_{j+1}(\x)$.
\end{itemize}

We emphasize on the fact again that we do not exactly compute $\widetilde{T}_{i,j+1}$ mod $z^{d_{j+1}}$; instead it is computed as a fraction in $\F(z,\x)$, with formal $z$. Formally, we compute $T_{i,j+1} \in \F(z,\x)$, such that $\widetilde{T}_{i,j+1} = T_{i,j+1}$, over $\ringR_{j+1}(\x)$. We keep track of the degree of $z$ in $T_{i,j+1}$. Next, we will show that all the inductive hypotheses assumed hold in the $j^{\text{th}}$ step as well. 

\smallskip

\noindent \textbf{Hypothesis (\ref{hypo1}): Definability of $T_{i,j+1}$ and $f_{j+1}$.} By the minimal valuation assumption, it follows that $\val(f_j) \ge v_{k-j,j}$, and thus $\widetilde{T}_{i,j+1}$ and $f_{j+1}$ are all well-defined over $\ringR_{j+1}(\x)$. Note that, \autoref{eq:induction-thm1-main-eq} holds over $\ringR_{j+1}(\x)$ as $d_{j+1} < d_j$ (because, whatever identity holds true $\bmod \ z^{d_j}$ must hold $\bmod \ z^{d_{j+1}}$ as well). Hence, we must have $\sum_{i=1}^{k-j-1}\, \widetilde{T}_{i,j+1} = f_{j+1}$, over $\ringR_{j+1}(\x)$ thus proving the induction hypothesis (\ref{hypo1}). 
 
%

\smallskip

\noindent \textbf{Hypothesis (\ref{hypo2}): Positivity of Valuation.} Since we divide by the $\min \val$, by definition we immediately get $\val_{z}(T_{i,j+1}) \ge 0$ proving the hypothesis. Further, we claim that min $\val$ computation in $\DiDI$ is easy. For this, recall from the definition of valuation \[\min_{i}\val_{z}(P_{i,j}/Q_{i,j})= \min_{i} (\val_{z}(P_{i,j}) - \val_{z}(P_{i,j}) ). \] Therefore, for $\min \val$ we compute $\val_{z}(P_{i,j})$ and $\val_{z}(Q_{i,j})$ for all $i \in [k-j]$.

Here is an important lemma which shows that coefficient of $y^e$ of a polynomial $f(\x,y) \in \F[\x,y]$, computed by a $\SWSU{}{}$ circuit, can be computed by a small $\SWSU{}{}$ circuit.

\begin{lemma}[Coefficient extraction]\label{lem:cf-swsw}
Let $f(\x,y) \in \F[y][\x]$ be computed by a $\SWSU{}{}$ circuit of size $s$ and degree $d$. Then, $\cf_{y^e}(f) \in \F[\x]$ can be computed by a small $\SWSU{}{}$ circuit of size $O(sd)$, over $\F[\x]$.
\end{lemma}
\begin{proof-sketch}
    Let, $f= \sum_{i} \alpha_i \cdot g_i^{e_i}$. Of course, $e_i \le s$ and $\pdeg_{y}(f) \le d$. Thus, write $f=\sum_{i=0}^{d}\,f_i\cdot y^i$, where $f_i \in \F[\x]$. We can interpolate on $d+1$-many distinct points $y\in \F$ and conclude that $f_i$ has a $\SWSU{}{}{}$ circuit of size at most $O(sd)$. 
\end{proof-sketch}

\noindent Using \autoref{lem:cf-swsw} we known $\cf_{z^e}(P_{i,j})$ and $\cf_{z^e}(Q_{i,j})$ are in $\SWSU{}{}$ over $F[\x]$. We can keep track of $z$ degree and thus interpolate to find the minimum $e < d_j$ such that the computed coefficients are $\ne 0$, which gives the respective $\val$.

\smallskip

\noindent \textbf{Hypothesis (\ref{hypo3}): The `iff' condition.} The above \autoref{eq:induction-thm1-main-eq} pioneers to reduce from $k-j$-summands to $k-j-1$. But we want a $\iff$ condition to efficiently reduce the identity testing. If $f_{j+1} \ne 0$, then $\val_{z}(f_{j+1}) < d_{j+1}$.  Further, $f_{j+1} = 0$, over $\ringR_{j+1}(\x)$ implies--
\begin{enumerate}
    \item Either, $f_j/T_{k-j,j}$ is $z$-free. This implies it is in $\F(\x)$. Now, if indeed $f_0 \ne 0$, then the computed $T_{i,j}$ as well as $f_j$ must be non-zero over $\F(z,\x)$, by induction hypothesis (as they are non-zero over $\ringR_j(\x)$).  However,
    \begin{align*}
    \left(\frac{T_{i,j}}{T_{k-j,j}}\right) \bigg\rvert_{z=0} &= \left( \frac{U_{i,j} \cdot V_{k-j,j}}{U_{k-j,j} \cdot V_{i,j}} \right) \bigg\rvert_{z=0} \cdot \left( \frac{P_{i,j} \cdot Q_{k-j,j}}{P_{k-j,j} \cdot Q_{i,j}} \right) \bigg\rvert_{z=0} \\ & \in\; \F \cdot \left( \frac{\SWSU{}{}{}}{\SWSU{}{}{}} \right).
    \end{align*} Thus, \[ \frac{f_j}{T_{k-j,j}} \; \in \; \sum \;  \F \cdot \left( \frac{\SWSU{}{}{}}{\SWSU{}{}{}} \right) \,\in\,  \left( \frac{\SWSU{}{}{}}{\SWSU{}{}{}} \right).\] Here we crucially use that $\SWSU{}{}{}$ is closed under multiplication (\autoref{lem:prod-swsw}). 
%
    Thus, this identity testing can be done in $\poly$-time (\autoref{lem:swsu-hs}).  For, detailed time-complexity and calculations, see \autoref{cl:size-bound-on-Ti} and its subsequent paragraph. \smallskip 
    \item Or, $\partial_{z}(f_j/T_{k-j,j}) = z^{d_{j+1}} \cdot p$, where $p \in \F(z,\x)$ s.t.~$\val_{z}(p) \ge 0$. By a simple power series expansion, one concludes that $p \in \F(\x)[[z]]$ (\autoref{lem:val-powerseries}). Hence, one concludes that \[ \frac{f_j}{T_{k-j,j}} \in  \left\langle z^{d_{j+1}+1}\right\rangle_{\F(\x)[[z]]}\; \implies \val_{z}(f_j) \ge d_j,\] i.e.~$f_j = 0$, over $\ringR_j(\x)$.
\end{enumerate}

\medskip
Conversely, $f_j=0$, over $\ringR_j(\x)$, implies \begin{align*}\val_{z}(f_j) \ge d_j &\implies \val_{z} \left( \partial_{z}\left( \frac{f_j}{T_{k-j,j}} \right) \right) \ge d_j-v_{k-j,j}-1 \\ &\implies f_{j+1}=0, \; \text{over} \; \ringR_{j+1}(\x).\end{align*}
Thus, we have proved that $\sum_{i \in [k-j]}\,T_{i,j}\,\ne 0 \;\text{ over}\;\ringR_j(\x)$ iff \[\sum_{i \in [k-j-1]}\, T_{i,j+1} \ne 0 \;\text{ over} \;\ringR_{j+1}(\x)\;,\,\text{or}\,,\;\;0 \ne \left( \frac{f_j}{T_{k-j,j}} \right)\bigg\rvert_{z=0} \in \F(\x)\;.\]
Therefore induction hypothesis (\ref{hypo3}) holds. 

%

\smallskip
\noindent \textbf{Hypothesis (\ref{hypo4}): Size analysis.} We will show that $T_{i,j+1} \in (\PSU{}{}/\PSU{}{}) \cdot (\SWSU{}{}/\SWSU{}{})$, over $\ringR_{j+1}(\x)$, with only polynomial blowup in size.  Let $\size(T_{i,j}) \le s_j$, for $i \in [k-j]$, and $j \in [k]$. Note that, by assumption, $s_0 \le s$.

\begin{claim}[Final size] \label{cl:size-bound-on-Ti}
$T_{1,k-1} \in (\PSU{}{}/\PSU{}{}) \cdot (\SWSU{}{}/\SWSU{}{})$ of size $s^{O(k7^k)}$, over $\ringR_{k-1}(\x)$.
\end{claim}

\begin{proof}
Steps $j=0$ and $j>0$ are slightly different because of the $\Phi$. However the main idea of using power-series is the same which eventually shows that $\dlg(\Sigma \wedge) \in \SWSU{}{}{}$. 

We first deal with $j=0$. Let $A-z \cdot B= \Phi(g) \in \Sigma \wedge$, for some $A \in \F$ and $B \in \ringR_1[\x]$. Note that $A \ne 0$ because of the map $\Psi$. Further, $\size(B) \le O(d \cdot \size(g))$, as a single monomial of the form $x^e$ can produce $d+1$-many monomials. Over $\ringR_1(\x)$,
\begin{align}
    \dlg(\Phi(g)) = -\frac{\partial_{z}(B\cdot z)}{A(1- \frac{B}{A}\cdot z)}  = -\frac{\partial_{z}(B\cdot z)}{A} \cdot \sum_{i=0}^{d_1-1} \left(\frac{B}{A}\right)^{i} \cdot z^i\;. \label{powerseries-1/g-phi}
\end{align}
$B^i$ has a trivial $\WSU{}$-circuit of size $O(d \cdot \size(g))$.  Also, $\partial_{z}(B\cdot z)$ has a $\Sigma \wedge$-circuit of size at most $O(d\cdot \size(g))$. Using waring identity (\autoref{lem:waring-rank}), we get that each $\partial_{z}(B\cdot z) \cdot (B/A)^i \cdot z^i$ has size $O(i\cdot d \cdot \size(g))$, over $\ringR_1(\x)$. Summing over $i \in [d_1-1]$, the overall size is at most $ O(d_1^2\cdot d \cdot \size(g))=O(d^3 \cdot \size(g))$, as $d_0=d_1=d$. 

For the $j$-th step, we emphasize that the degree could be larger than $d$. Assume that syntactic degree of denominator and numerator of $T_{i,j}$ (each in $\F[\x,\z]$) are bounded by $D_{j}$ (it is {\em not} $d_j$ as seen above; this is to save on the trouble of mod-computation at each step). Of course, $D_0<d \le s$. 

For $j>0$, the above summation in \autoref{powerseries-1/g-phi} is over $\ringR_{j}(\x)$. However the degree could be $D_j$ (possibly more than $d_j$) of the corresponding $A$ and $B$. Thus, the overall size after the power-series expansion would be $O(D_j^2 \cdot d \cdot \size(g))$. 

Using \autoref{lem:diff-size}, we can show that $\dlg(P_{i,j}) \in \SWSU{}{}/\SWSU{}{}$ (similarly for $Q_{i,j}$), of size $O(D_{j}^2 \cdot s_j)$. Also $\dlg(U_{i,j} \cdot V_{k-j,j}) \in \sum \,\dlg(\Sigma \wedge)$, i.e. sum of action of $\dlg$ on $\Sigma \wedge$ (since $\dlg$ linearizes product); and it can be computed by the above formulation. 
Thus, $\dlg(T_{i,j}/T_{k-j,j})$ is a sum of $4$-many  $\SWSU{}{}/\SWSU{}{}$ of size at most $O(D_{j}^2\,s_j)$ and $1$-many $\SWSU{}{}$ of size $O(D_j^2 d_j s_j)$ (from the above power-series computation) [Note: we summed up the $\SWSU{}{}$-expressions from $\dlg(\Sigma \wedge)$ together]. Additionally the syntactic degree of each denominator and numerator (of the $\SWSU{}{}{}/\SWSU{}{}{}$) is $O(D_j)$. We rewrite the $4$ expressions (each of $\SWSU{}{}{}/\SWSU{}{}{}$) and express it as a single $\SWSU{}{}/\SWSU{}{}$ using waring identity (\autoref{lem:prod-swsw}), with the size blowup of $O(D_{j}^{12}\,s_j^{4})$; here the syntatic degree blowsup to $O(D_j)$. Finally we add the remaining $\SWSU{}{}$ circuit (of size $O(D_j^3 s_j)$ and degree $O(dD_j)$) to get $O(s_j^{5} D_{j}^{16}d)$. To bound this, we need to understand the degree bound $D_j$. 

Finally we need to multiply $T_{i,j}/T_{k-j,j} \in (\PSU{}{}/\PSU{}{}) \cdot (\SWSU{}{}/\SWSU{}{})$ where each $\SWSU{}{}$ is a product of two $\SWSU{}{}$ expression of size $s_j$ and syntactic degree $D_j$; clubbed together owing a blowup of $O(D_j \cdot s_j^2)$. Hence multiplying it with $\SWSU{}{} / \SWSU{}{}$ expression obtained from $\dlg$ computation above gives size blowup of $s_{j+1} = s^7 \cdot D_j^{O(1)} \cdot d$.

Computing $T_{i,j}/T_{k-j,j}$ increases the syntactic degree `slowly'; which is much less than the size blowup.
As mentioned before, the deg-blowup in $\dlg$-computation is $O(dD_j)$ and in the clearing of four expressions, it is just $O(D_j)$. Thus, $D_{j+1}= O(dD_j) \implies D_j= d^{O(j)}$. 
 
The recursion on the size is $s_{j+1} = s_j^{7} \cdot d^{O(j)}$. Using $d \le s$ we deduce, $s_j= (sd)^{O(j \cdot 7^j)}$. In particular, $s_{k-1}$, size after $k-1$ steps is $s^{O(k \cdot 7^k)}$. This computation quantitatively establishes induction hypothesis (\ref{hypo4}).
\label{deg-bound-thm1}
\end{proof}

\smallskip

\noindent \textbf{Hypothesis (\ref{hypo5}): Invertibility of $\PSU{}{}$-circuits.} For invertibility, we want to emphasise that the $\dlg$ compuation plays a crucial role here. In the following lemma we claim that the action $\dlg(\SWSU{}{}) \in \SWSU{}{}{}/\SWSU{}{}{}$, is of $\poly$-size. 

\begin{lemma}[Differentiation] \label{lem:diff-size}
    Let $f(\x,y) \in \F[y][\x]$ be computed by a $\SWSU{}{}$ circuit of size $s$ and degree $d$. Then, $\partial_{y}(f)$ can be computed by a small $\SWSU{}{}$ circuit of size $O(sd^2)$, over $\F[y][\x]$.
\end{lemma}
\begin{proof-sketch}
    \autoref{lem:cf-swsw} shows that each $f_e$ has $O(sd)$ size circuit where $f=\sum_{e} f_e\,y^e$. 
    Doing this for each $e \in [0,d]$ gives a blowup of $O(sd^2)$.
\end{proof-sketch}

Similarly consider the action on $ \PSU{}{}$. We know $\dlg$ distributes the product additively, so it suffices to work with $\dlg(\Sigma \wedge)$; and earlier in \autoref{cl:size-bound-on-Ti} we saw that $\dlg(\Sigma \wedge) \in \SWSU{}{}{}$ of $\poly$-size. Assuming these, we simplify \[ \frac{T_{i,j}}{T_{k-j,j}} = \frac{U_{i,j} \cdot V_{k-j,j}}{V_{i,j} \cdot U_{k-j,j}} \cdot \frac{P_{i,j} \cdot Q_{k-j,j}}{Q_{i,j} \cdot P_{k-j,j}},\] and its $\dlg$. Thus, using \autoref{eq:induction-thm1-main-eq}, $U_{i,(j+1)}$ grows to $U_{i,j} \cdot V_{k-j,j}$ (and similarly $V_{i,(j+1)}$). This also means: $U_{i,(j+1)} \rvert_{z=0}$ $\in \F\setminus\{0\}$ and thereby proving the hypothesis.


\subsubsection*{Final time complexity} 
The above proof actually shows that $T_{1,k-1}$ is in $\Gen{1}{s^{O(k\cdot 7^k)}}$ over $\ringR_{k-1}(\x)$; and that the degree bound on $z$ (over $\F[z,\x]$, keeping denominator and numerator `in place') is $D_{k-1} = d^{O(k)}$. We cannot directly use the identity testing algorithms of the constituent simpler models due to $\ringR_{k-1}(\x)$. Moreover, using hypothesis (\ref{hypo2}) and \autoref{lem:val-powerseries} we know that $T_{1,k-1} \in \F(\x)[[z]]$ and it suffices to do identity testing on the first term of the powerseries: $T_{1,k-1} \rvert_{z = 0}$ over $\F(\x)$. Note that, hypothesis (\ref{hypo5}) guarantees that $\PSU{}{}$ part remains non-zero on $z = 0$ evaluation, however, $\SWSU{}{}/\SWSU{}{}$ may be undefined. For this, we keep track of $z$ degree of numerator and denominator, which will be polynomially bounded as seen in the discussion above. We can easily interpolate and cancel the $z$ power to make it work. Basically this shows that to test $T_{1,k-1}$ we need to test $z^e \cdot \SWSU{}{}$ over $\F[\x]$ where $e \geq 0$ due to positive valuation. Whitebox PIT of $\SWSU{}{}$ is in poly-time using \autoref{lem:swsu-hs}, and testing $z^e$ is possible using \autoref{lem:trivial-pit} with appropriate degree bound. The proof above is constructive: we calculate $U_{i,j+1}$ (and other terms) from $U_{i,j}$ explicitly. Gluing everything together we conclude this part can be done in $s^{O(k 7^k)}$ time.

What remains is to test the $z=0$-part of induction hypothesis (\ref{hypo3}); it could \emph{short-circuit} the recursion much before $j=k-1$. As we mentioned before, in this case, we need to do a PIT on $\SWSU{}{}{}$ only. At the $j$-th step, when we substitute $z=0$, the size of each $T_{i,j}$ can be at most $s_j$ (by definition). We need to do PIT on a simpler model: $\sum^{[k-j]}\;\F \cdot (\SWSU{}{}{}/\SWSU{}{}{})$. We can clear out and express this as a single $\SWSU{}{}{}/\SWSU{}{}{}$ expression; with a size blowup of $s_j^{O(k-j)} \le (sd)^{O(j(k-j)7^j)}$. Since this case could short-circuit the recursion, to bound the final time complexity, we need to consider the $j$ which maximizes the exponent.

\begin{lemma}\label{max-exp-opt}
    Let $k \in \N$, and $h(x):= x (k-x)7^x$. Then, $\max_{i \in [k-1]} h(i)=h(k-1)$.
\end{lemma}
\begin{proof-sketch}
    Differentiate to get $h'(x)= (k-x)7^x - x7^x +  x(k-x)(\log 7)7^x = 7^x\cdot[x^2(-\log7) + x(k\log7-2) + k ]$. It vanishes at \[x = \left(\frac{k}{2} - \frac{1}{\log 7}\right) + \sqrt{\left(\frac{k}{2} - \frac{1}{\log 7}\right)^2 - \frac{k}{\log 7}}\;.\]Thus, $h$ is maximized at the integer $x=k-1$.
\end{proof-sketch}


Therefore, $\max_{j \in [k-1]} j(k-j)7^j = (k-1)7^{k-1}$. 
%
Finally, use \autoref{lem:swsu-hs} for the base-case whitebox PIT. Thus, the final time complexity is $s^{O(k \cdot 7^k)}$. 

Here we also remark that in $z =0$ substitution $\SWSU{}{}/\SWSU{}{}$ may be undefined. However, we keep track of $z$ degree of numerator and denominator, which will be polynomially bounded as seen in the discussion above. We can easily interpolate and cancel the $z$ power to make it work. 

\smallskip
{\bf \noindent Bit complexity.} It is routine to show that the  bit-complexity is really what we claim. Initially, the given circuit has bit-complexity $s$. The main blowup happens due to the $\dlg$-computation which is a poly-size blowup. We also remark that while using \autoref{lem:prod-swsw} (using \autoref{lem:waring-rank}), we \emph{may} need to go to a field extension of at most $s^{O(k)}$ (because of the $\varepsilon(i)$ and correspondingly the constants $\gamma_{\varepsilon(2),\hdots,\varepsilon(k)}$, but they still are $s^{O(k)}$-bits). Also, \autoref{thm:sparse-ks} and \autoref{lem:swsu-hs} computations blowup bit-complexity polynomially. This concludes the proof.
\end{proof}

\begin{remark*}
\begin{enumerate}
\setlength\itemsep{.1em}
    \item The above method does {\em not} give whitebox PIT (in $\poly$-time) for $\SPSP{[k]}{}{[\delta]}$, as we donot know $\poly$-time whitebox PIT for $\SWSP{}{}{[\delta]}$. However, the above methods do show that whitebox-PIT for $\SPSP{[k]}{}{[\delta]}$ polynomially \emph{reduces} to whitebox-PIT for $\SWSP{}{}{[\delta]}$.
    \item DiDI-technique can be used to give whitebox PIT for the general bloated model $\Gen{k}{s}$. 
    \item The above proof works when the characteristic is $\ge d$. This is because the nonzeroness remains \emph{preserved} after derivation wrt $z$.
\end{enumerate}
\end{remark*}

\subsection{Algorithm}

The whitebox PIT for \autoref{thm:thm1}, that is discussed in section \ref{sec:pf:thm1}, appears (below) as Algorithm \ref{algo:whitebox-algo}.

\begin{algorithm}
    \caption{Whitebox PIT Algorithm for $\SPSU{[k]}{}{}$-circuits \label{algo:whitebox-algo}}
    \hspace*{\algorithmicindent} \textbf{INPUT: $f=T_{1} + \hdots + T_k \in \SPSU{[k]}{}{}$, a whitebox circuit of size $s$ over $\F[\x]$} \\
    \hspace*{\algorithmicindent} \textbf{OUTPUT: $ 0$, if $f \equiv 0$, and $1$, if non-zero.} \\
%
    \begin{algorithmic}[1]
    \STATE{Let $\Psi : \F[\x] \longrightarrow \F[z]$, be a sparse-PIT map, using \cite{klivans2001randomness} (\autoref{thm:sparse-ks}). Apply it on $f$ and  check  whether $\Psi(f) \stackrel{?}{=} 0$. If non-zero, {\tt output} $1$}

    \STATE{Obtain a point $\alphavec = (a_1, \dots, a_n) \in \F^n$ from Hitting Set $\calH$ of $\PSU{}{}$ such that $T_{i} \rvert_{\x = \alphavec} \neq 0$, for all $i \in [k]$. And define $\Phi : x_i \mapsto z \cdot x_i + a_i$. Check $\sum_{i \in [k-1]} \partial_{z}(\Phi(T_i)/\Phi(T_k)) \stackrel{?}{=} 0 \bmod z^{d_1}$ ($d_1:=s$) as follows:}
%
    
    \STATE{Consider each $T_{i,1}:=\partial_{z}(\Phi(T_i)/\Phi(T_k))$ over $R_1(\x)$, where $R_1:=\F[z]/\langle z^{d_1} \rangle$. Use $\dlg$ computation (\autoref{cl:size-bound-on-Ti}), to write each $T_{i,1}$ in a `bloated' form as $(\PSU{}{}/\PSU{}{}) \cdot (\SWSU{}{}{}/\SWSU{}{}{})$.}
    
    \FOR{$j\leftarrow 1$ \TO $k-1$}
        \STATE{Reduce the top-fanin at each step using `{\bf Di}vide \& {\bf D}erive' technique. Assume that at $j$-th step, we have to check the identity:}
        $\sum_{i \in [k-j]} T_{i,j}\; \stackrel{?}{=}\;0\; \text{over}\;R_j(\x),\;\text{where}\;R_j:= \F[z]/\langle z^{d_j}\rangle\;$, 
        each $T_{i,j}$ has a $(\PSU{}{}/\PSU{}{}) \cdot (\SWSU{}{}{}/\SWSU{}{}{})$ representation and therein each $\PSU{}{} \rvert_{z=0} \in \F\setminus\{0\}$. 
        
        \STATE{Compute $v_{k-j,j}:= \min_i \val_{z}(T_{i,j})$; by reordering it is for $i=k-j$. 
        To compute $v_{k-j,j}$, use coefficient extraction (\autoref{lem:cf-swsw}) and $\SWSU{}{}{}$-circuit PIT (\autoref{lem:swsu-hs}).}
            
        \STATE{`{\bf D}ivide' by $T_{k-j,j}$ and check whether $\left(\sum_{i \in [k-j-1]}\; (T_{i,j} / T_{k-j,j}) +1 \right)\bigg\rvert_{z=0} \stackrel{?}{=} 0$. Note: this expression is in $(\SWSU{}{}{}/\SWSU{}{}{})$. Use--- (1) $\PSU{}{}\rvert_{z=0} \in \F$, and (2) \emph{closure} of $\SWSU{}{}{}$ under multiplication. Finally, do PIT on this by \autoref{lem:swsu-hs}.}
            
        \STATE{If it is non-zero, {\tt output} $1$, otherwise `{\bf D}erive' wrt $z$ and `{\bf I}nduct' on $\left(\sum_{i \in [k-j-1]} \partial_{z}(T_{i,j}/T_{k-j,j})\right) \stackrel{?}{=} 0$, over $R_{j+1}(\x)$ where $R_{j+1}:= \F[z]/\langle z^{d_j-v_{k-j,j}-1} \rangle$.}
            
        \STATE{Again using $\dlg$ (\autoref{cl:size-bound-on-Ti}), show that $T_{i,j+1}:=\partial_{z}(T_{i,j}/T_{k-j,j})$ has small $(\PSU{}{}/\PSU{}{}) \cdot (\SWSU{}{}{}/\SWSU{}{}{})$-circuit over $R_{j+1}(\x)$. So call the algorithm on $\sum_{i \in [k-j-1]}\, T_{i,j+1} \stackrel{?}{=}0$.}
        
        \STATE{$j \leftarrow j+1$.}
    \ENDFOR
    
    \STATE{At the end, $j=k-1$. Do PIT (\autoref{lem:swsu-hs}) on the single $(\PSU{}{}/\PSU{}{})\cdot (\SWSU{}{}{}/\SWSU{}{}{})$ circuit, over $R_{k-1}(\x)$. If it is zero, {\tt output} $0$ otherwise {\tt output} $1$.}
    \end{algorithmic}
    \end{algorithm}
    
\noindent \emph{Words of caution}: Throughout the algorithm there are intermediate expressions to be stored compactly. Think of them as `special' circuits in $\x$, but over the {\em function-field} $\F(\z)$. Keep track of their degrees wrt $z$;  and that of the sizes of their fractions represented in `bloated' circuit form. 
